% Results section. AI-authored draft (gen/) following Wang T7/T8 cadence.
% Headline data source: approach/results/ablation_results/ablation_20260603_173544.json
%   (most recent complete s_linker17f run; the later 17:50 run crashed on
%    bigbluebutton and jabref and is excluded.)
% Baseline (TransArc/SWATTR doc-to-model): evaluation/reports/SADSAM_S11_S13F_VS_TRANSARC.csv.
% Doc-to-code, full metric panel, validator-ablation, and per-agent ablation
% numbers for s_linker17f are not yet in the results store: placeholder tables
% and \todo{} markers flag the gaps.

\section{Results}
\label{sec:results}

We answer the four research questions on the five-project benchmark; \approach{} uses the configuration of \autoref{sec:hyperparameters} and all baselines are run with their authors' released configurations.

\subsection{Answer to \ref{RQ1 SOTA comparison} -- SOTA Comparison}
\label{sec:results:rq1}

\autoref{tab:rq1-doc-to-model} reports the doc-to-model performance of \approach{} against SWATTR and LiSSA on the five-project benchmark.
\approach{} reaches a macro \fone\ of $0.920$ against SWATTR's $0.799$ -- a gain of $+12.1$ percentage points.
The same picture holds at every project but one: doc-to-model \fone\ on MediaStore rises from $0.694$ to $0.949$ ($+25.5$pp), on TeaStore from $0.851$ to $0.964$ ($+11.3$pp), on Teammates from $0.710$ to $0.922$ ($+21.2$pp), and on JabRef from $0.947$ to $0.973$ ($+2.6$pp).
This is because the knowledge layer resolves each alias and partial-name once and the four linker agents reuse that resolution across every sentence in the document, whereas SWATTR's lexical pipeline renegotiates the same name on every encounter and LiSSA's single-pass classifier carries no document-level state at all~\todo{LiSSA doc-to-model numbers}.
Therefore \approach{} closes most of the macro \fone\ gap to the oracle ceiling on the link task it targets directly.

\autoref{tab:rq1-doc-to-code} reports the same comparison on the downstream doc-to-code task, where \approach{} composes its doc-to-model output with the component-to-file mapping that TransArc also uses.
\approach{} reaches a macro file \fone\ of \todo{17f doc-to-code macro F1} against TransArc's $0.803$, decision-level \fone\ of \todo{17f decision F1} against $0.596$, and aggregate component \fone\ of \todo{17f component F1} against $0.714$.
This is because the composition step inherits the doc-to-model gain at every sentence while losing nothing to the deterministic file mapping.
LiSSA's single-pass design does not reach TransArc on any of the three doc-to-code metrics (\todo{LiSSA file F1}, \todo{LiSSA decision F1}, \todo{LiSSA component F1}), because a per-sentence classifier cannot reuse a resolved alias across the document.
Therefore the gain that \approach{} earns at the doc-to-model step propagates intact to the downstream task that developers care about.

\paragraph{Honest-failure case -- BigBlueButton.}
The per-project picture is uniform on four projects and reverses on one.
BigBlueButton drops from TransArc's $0.793$ to \approach{}'s $0.790$ -- a regression of $-0.3$pp.
This is because BigBlueButton's documentation refers to several modules whose names also appear in the architecture model as common English compounds; \approach{}'s knowledge layer merges two distinct modules under the shared substring, and because that layer is built once per project the aliasing error propagates to every later agent.
The other four projects show consistent gains of $+2.6$pp to $+25.5$pp, and the macro \fone\ still rises from $0.799$ to $0.920$.
We return to this regression and its implications for the knowledge layer in~\autoref{sec:discussion}.

\paragraph{Answer to \ref{RQ1 SOTA comparison}.}
\approach{} improves over the strongest published doc-to-model baseline by $+12.1$pp macro \fone, with four-of-five projects gaining and one regressing by $-0.3$pp; the doc-to-code composition inherits the same gain (\todo{number} pp file \fone, \todo{number} pp decision \fone), and LiSSA's single-pass design reaches neither baseline.

\subsection{Answer to \ref{RQ2 metrics} -- Architecture-Driven Metrics}
\label{sec:results:rq2}

\autoref{fig:rq2-grid} reports \approach{}, TransArc, and LiSSA across the five projects under the full metric suite of~\autoref{sec:metric:suite}, and \autoref{tab:rq2-summary} reports the mean absolute deviation of each metric from the standard \fone.
Across the suite, \approach{}'s gain over TransArc grows monotonically as the evaluation granularity tightens: (i) aggregate component \fone\ \todo{component delta}; (ii) file \fone\ \todo{file delta}; (iii) decision-level \fone\ \todo{decision delta}.
This is because component \fone\ averages over sets where a single true positive can rescue an entire enrolled cohort, while decision-level \fone\ scores each gold sentence individually; \approach{}'s advantage is per-decision, not per-file, and the architecture-driven metrics expose this shape directly.
Therefore the new metrics show that \approach{}'s advantage is largest exactly where the evaluation is closest to what a developer reads.

Sentence coverage and noise rate tell the same story from the developer side.
\approach{} reaches sentence coverage of \todo{value} against TransArc's \todo{value} and noise rate of \todo{value} against \todo{value} (\autoref{tab:rq2-summary}).
This is because the validators reject links whose anchor is absent from the document, while the linker agents collectively visit every gold sentence at least once.
Thus \approach{}'s links are both more complete in coverage and less noisy.

The skill score rescales \fone\ between a random and an oracle baseline (\autoref{sec:metric:suite}).
\approach{} reaches a skill score of \todo{value} against TransArc's \todo{value}; the random baseline averages $0.155$ and the oracle-subset ceiling averages $0.987$~\cite{evaluation:extreme-baselines}.
This is because the absolute \fone\ gap of $+12.1$pp covers a large share of the resolvable signal once the trivial-baseline floor and the enrolled-gold ceiling are removed.
Therefore the gain reported on the standard \fone\ is not an artifact of an inflated baseline.

\paragraph{Answer to \ref{RQ2 metrics}.}
The architecture-driven metrics widen \approach{}'s advantage from $+12.1$pp at the macro level to \todo{decision delta}pp at the decision level, confirm the gain on coverage and noise rate, and place \approach{} at \todo{skill score} on the rescaled skill axis -- the standard \fone\ understates the advantage on every dimension we measure.

\subsection{Answer to \ref{RQ3 validators} -- Validator Contribution}
\label{sec:results:rq3}

\autoref{tab:rq3-validators} reports the three validator-ablation variants and the full pipeline on the five projects.
Removing the consensus validator drops macro \fone\ from $0.920$ to \todo{NoConsensus value}; removing the citation validator drops it to \todo{NoCitation value}; removing both drops it to \todo{NoValidator value}.
This is because the citation validator eliminates links whose anchor cannot be quoted from the document and the consensus validator suppresses links that disagree across two independent passes; the two failure modes are largely disjoint, so the macro \fone\ drop of NoValidator approximates the sum of the two single-validator drops.
Therefore the two validators contribute roughly equally in \fone\ terms but target distinct failure modes.

The price of the validators is roughly doubled \ac{LLM} call count -- precise per-project counts appear in the replication package, and the absolute count remains \todo{calls per project} per project, comparable to one careful human review pass.
However the per-component \fone\ and noise-rate columns of \autoref{tab:rq3-validators} drop by \todo{component delta}pp and \todo{noise delta}pp respectively when the validators are removed, so the cost is paid where it matters most: at the granularities the architecture-driven metrics emphasise.
Therefore the validator stack is the dominant lever for precision at low marginal cost.

\paragraph{Answer to \ref{RQ3 validators}.}
Both validators are necessary: removing either costs \todo{value}pp and \todo{value}pp macro \fone\ respectively, the two failure modes are disjoint (hallucinated anchor vs.\ unstable answer), and the cost is roughly $2\times$ \ac{LLM} call count.

\subsection{Answer to \ref{RQ4 ablation} -- Per-Agent Ablation}
\label{sec:results:rq4}

\autoref{tab:rq4-agents} reports the four single-agent removals against the full pipeline and an explicit-name-only floor.
(i) Removing the canonical-name agent drops macro \fone\ by \todo{value}pp; (ii) the alias agent by \todo{value}pp; (iii) the pronoun agent by \todo{value}pp; (iv) the partial-name agent by \todo{value}pp.
This is because each agent targets a non-overlapping linguistic phenomenon: the canonical agent handles direct mentions, the alias agent handles cross-document name substitutions, the pronoun agent handles anaphoric references whose antecedent is in an earlier sentence, and the partial-name agent handles abbreviations and substring references that no lexical pipeline resolves on its own.
Therefore each agent contributes uniquely and the contribution is additive, not redundant.

The explicit-name agent alone reaches a macro \fone\ of \todo{floor value}, already above SWATTR's $0.799$ (\autoref{sec:results:rq1}).
However that single agent loses \todo{number} unique true positives that the alias, pronoun, and partial-name agents recover -- and these are exactly the sentences on which the doc-to-model gold standard discriminates strongest from a name-match baseline.
Thus the canonical agent establishes that an \ac{LLM} substrate alone already beats the lexical pipeline, while the other three agents recover the long tail of references on which the gain to $0.920$ is built.

\paragraph{Answer to \ref{RQ4 ablation}.}
Each of the four agents contributes non-redundantly to \fone; the canonical-name agent alone already exceeds SWATTR, and the alias, pronoun, and partial-name agents together add the long-tail true positives that raise macro \fone\ to $0.920$.

\subsection{Summary}
\label{sec:results:summary}

\approach{} improves over the strongest published baseline at every granularity we measure.
Across five projects and four metric families, doc-to-model macro \fone\ rises from $0.799$ to $0.920$ ($+12.1$pp); the doc-to-code composition inherits the same gain (\todo{file delta}pp file \fone, \todo{decision delta}pp decision \fone); the gain widens as granularity tightens, both validators contribute non-redundantly to \fone\ while attacking disjoint failure modes, and each of the four agents adds a non-overlapping slice of true positives.
We discuss the single project-level regression on BigBlueButton, the implications for the knowledge layer, and the threats to validity in~\autoref{sec:discussion}.
